\documentclass[10pt,twocolumn]{article}
\usepackage[margin=0.75in]{geometry}
\usepackage{mathptmx}
\usepackage[T1]{fontenc}
\usepackage{microtype}
\usepackage{booktabs}
\usepackage{enumitem}
\usepackage[hidelinks]{hyperref}
\setlist{nosep,leftmargin=*}

% Every re-analysis number below is in paper/reanalysis.json, produced by
% paper/reanalysis.py. The recorded 2025 arms come from
% data/experiments/final_results_summary.md (also parsed into the JSON).

\title{When Constants Win: Lessons from Re-analysing\\a Privacy-Harm Heuristics Evaluation}
\author{Anonymous Author(s)}
\date{}

\begin{document}
\maketitle

\begin{abstract}
We re-analyse a small study that asked whether interpretable heuristics can
estimate privacy harm, and how they compare with large language models (LLMs).
Its recorded result was that a ``rules static'' arm beat a plain LLM, a
rules-dynamic arm, retrieval augmentation, and two hybrids on all four reported
metrics. Using only committed artifacts and no model calls, we find:
\begin{itemize}
  \item Two constants chosen in advance, predicting \emph{no harm} or
        \emph{all four} of Solove's harm groups for every case, match or beat
        every arm on exact match, Jaccard, and nDCG@5.
  \item On micro-F1, the recorded arm leads, but no per-case outputs survive.
        A 2026 rerun of it is indistinguishable from the best constant unless
        out-of-taxonomy labels are removed.
  \item In the committed code, five of the six arms sent an identical prompt.
  \item The gold set was LLM-generated and mostly off-topic.
  \item Every stored multi-arm experiment was a dry run on a different gold set.
  \item Every training label came from a keyword fallback.
\end{itemize}
We distil eight evaluation lessons and one data-ethics lesson.
\end{abstract}

\section{Introduction}

Estimating the \emph{harm} of a privacy incident, not only detecting that one
occurred, matters to regulators, auditors, and engineers deciding what to fix
first. Solove's taxonomy organises privacy problems into information
collection, processing, dissemination, and invasion~\cite{solove2006}, and
Citron and Solove catalogue the injuries that follow~\cite{citron2022}. Recent
work applies the taxonomy to AI incidents~\cite{lee2024}. A natural question is
whether cheap, interpretable heuristics can apply that structure at scale, and
whether an LLM does better.

A 2025 master's practicum asked exactly that: \emph{can we estimate privacy harm
using heuristics, given facts parsed by a small language model}, and \emph{how
do non-language models compare with LLMs and with an LLM applying classical
privacy frameworks}? Its public artifact records a headline table in which a
``rules static'' arm led every metric. We re-examined every committed artifact
offline and asked what that evidence supports. The contributions are:
\begin{itemize}
  \item a re-analysis against pre-specified constant predictors, scored with the
        study's own metric code, with paired bootstrap intervals for every gap
        (\S\ref{sec:baselines});
  \item an audit of what the evaluation actually compared: its prompts, scorer
        revisions, gold set, dry runs, and label sources (\S\ref{sec:audit});
  \item nine lessons, each tied to its evidence (\S\ref{sec:lessons}).
\end{itemize}

\section{The original study}
\label{sec:study}

\paragraph{Pipeline.} Public incident reports (regulator and securities
filings, news, forums, Wikipedia) were crawled, labelled with Solove categories,
and featurised. Five interpretable model families were trained: a decision
tree; an Explainable Boosting Machine~\cite{lou2012,lou2013,nori2019}; Bayesian
rule lists~\cite{letham2015}; a sparse linear model; and a Chow--Liu Bayes
net~\cite{chow1968}. The labeller was to use GLiNER2~\cite{zaratiana2025}, a
successor to GLiNER~\cite{zaratiana2024}.

\paragraph{Evaluation.} Six arms were scored on a 50-case gold set (the
\emph{v3} set): a plain LLM baseline; \emph{rules static}, a fixed rule block in
the prompt; \emph{rules dynamic}; RAG; and two hybrids. All were served by
Gemini. The metrics were instance-averaged Jaccard, exact-match ratio (EMR),
micro-F1, and nDCG@5~\cite{zhang2014,jarvelin2002} over Solove's four groups.
The recorded table is dated 2025-11-06 and reports means. The summary says
there were five trials, but all six EMRs are multiples of $1/150$, and three are
not multiples of $1/250$. That is consistent with three trials of 50 cases. A
2026 rerun repeated rules-static once. Its results file records
\texttt{model\_name: null} and no resolved model, so the model it used is known
only from README prose. The repository also added an \emph{offline} arm that
calls no LLM and uses the corpus keyword labeller; we recompute it per case.

\begin{table*}[t]
\centering\small
\begin{tabular}{@{}lccccccc@{}}
\toprule
& \multicolumn{5}{c}{All 50 cases} & \multicolumn{2}{c}{15 labelled cases} \\
\cmidrule(lr){2-6}\cmidrule(l){7-8}
Arm & Jaccard & Jaccard$_{\emptyset=1}$ & EMR & micro-F1 & nDCG@5 & Jaccard & nDCG@5 \\
\midrule
\multicolumn{8}{@{}l}{\emph{Recorded 2025 means (no per-case outputs survive)}}\\
Rules static                 & 0.068 & -- & 0.667 & \textbf{0.391} & 0.090 & -- & -- \\
LLM baseline                 & 0.056 & -- & 0.620 & 0.285 & 0.072 & -- & -- \\
Rules dynamic                & 0.055 & -- & 0.593 & 0.281 & 0.073 & -- & -- \\
Hybrid, deterministic first  & 0.053 & -- & 0.593 & 0.275 & 0.068 & -- & -- \\
RAG                          & 0.047 & -- & 0.580 & 0.242 & 0.061 & -- & -- \\
Hybrid, LLM first            & 0.036 & -- & 0.600 & 0.195 & 0.049 & -- & -- \\
\midrule
\multicolumn{8}{@{}l}{\emph{Later runs (per-case outputs available)}}\\
Rules static, 2026, current scorer   & 0.105 & 0.465 & 0.360 & 0.336 & 0.176 & 0.351 & 0.587 \\
\quad same outputs, 2025 scorer      & 0.041 & 0.401 & 0.360 & 0.167 & 0.041 & 0.136 & 0.136 \\
\quad same outputs, filtered to groups (post hoc) & 0.158 & 0.518 & 0.400 & \emph{0.463} & 0.202 & 0.528 & 0.675 \\
Offline keyword arm (no LLM)         & 0.081 & 0.161 & 0.080 & 0.222 & 0.162 & 0.269 & 0.541 \\
\midrule
\multicolumn{8}{@{}l}{\emph{Constants, pre-specified}}\\
Always $\emptyset$                   & 0.000 & \textbf{0.700} & \textbf{0.700} & 0.000 & 0.000 & 0.000 & 0.000 \\
Always all four, taxonomy order      & \textbf{0.175} & 0.175 & 0.040 & 0.298 & \textbf{0.237} & \textbf{0.583} & \textbf{0.790} \\
\multicolumn{8}{@{}l}{\emph{Constants, best per column (selected on test labels)}}\\
Best of 65 constants                 & 0.192 & 0.700 & 0.700 & 0.356 & 0.282 & -- & -- \\
\midrule
Perfect predictor (ceiling)          & 0.300 & 1.000 & 1.000 & 1.000 & 0.300 & 1.000 & 1.000 \\
\bottomrule
\end{tabular}
\caption{All rows use the same 50-case gold set and the repository's metric
code, which reproduces the 2026 run exactly. That scorer gives empty-vs-empty
a Jaccard and nDCG of 0 but an EMR of 1, so a perfect predictor scores only
0.300 on Jaccard and nDCG. Jaccard$_{\emptyset=1}$ uses the other common
convention. The 65 constants are every ordered subset of the four groups.
Bold marks the best value per column among the recorded arms, the 2026
current-scorer run, the offline arm, and the pre-specified constants. It
excludes the post-hoc filtered row (italic where it would win), the
test-selected constants, and the ceiling.}
\label{tab:main}
\end{table*}

\section{Constants against the arms}
\label{sec:baselines}

\paragraph{An empty-majority gold set.} Of the 50 gold cases, 35 carry no harm
label. The other 15 carry 35 group labels (processing 13, invasion 11,
collection 7, dissemination 4). The scorer counts an empty prediction on an
empty case as an exact match (EMR 1) but as a Jaccard and nDCG of 0. EMR
therefore mostly measures how often an arm predicts nothing, and Jaccard and
nDCG are read against a ceiling of 0.30 (Table~\ref{tab:main}).

\paragraph{Pre-specified constants.} We fixed two label-free constants before
looking at scores: the empty set, and all four groups in the taxonomy's own
order.
\begin{itemize}
  \item \emph{Always empty} scores EMR 0.700. That is higher than every arm,
        although the 2025 rules-static EMR (0.667) is within about two cases
        and cannot be tested without per-case outputs.
  \item \emph{All four} beats every arm on Jaccard (0.175) and nDCG@5 (0.237),
        and so does each of its 24 orderings (nDCG 0.208--0.282).
  \item Paired case-resampling bootstrap intervals (2{,}000 resamples)
        against the 2026 rerun exclude zero: empty minus rerun EMR
        $[0.22, 0.46]$; all-four minus rerun Jaccard $[0.033, 0.110]$ and
        nDCG $[0.018, 0.111]$.
  \item Against the offline arm, all three of all-four's gaps exclude zero:
        Jaccard $[0.042, 0.152]$, micro-F1 $[0.010, 0.149]$, nDCG
        $[0.010, 0.146]$.
  \item Under the other Jaccard convention (empty-vs-empty = 1), the empty
        constant wins Jaccard as well. A constant wins either way; the
        convention decides which one.
\end{itemize}

\paragraph{Micro-F1, the exception.} The recorded rules-static micro-F1 (0.391)
exceeds even the best constant chosen with the test labels (0.356,
\{processing, invasion\}), but no per-case outputs survive to test it. The
2026 rerun scores 0.336. Against the best constant re-selected inside each
resample, the paired interval is $[-0.086, 0.108]$: no detectable difference.
The rerun's output contains eight labels outside the gold label space
(root-cause terms such as \texttt{lack\_of\_consent}), and each is a
guaranteed false positive. A post-hoc sensitivity check drops them. The
filtered output reaches micro-F1 0.463, and its interval against the
re-selected best constant is $[0.005, 0.255]$. So the arm carries signal that
its prompt's label vocabulary buries. Because the filter was chosen after we
saw the failure, this is a hypothesis for a new run, not a result.

\paragraph{The labelled cases alone.} On the 15 non-empty cases, the all-four
constant still leads on Jaccard (0.583 vs.\ 0.351 raw, 0.528 filtered) and
nDCG@5 (0.790 vs.\ 0.587 and 0.675). On micro-F1 it roughly ties the filtered
output (0.737 vs.\ 0.746). With only 15 informative cases, the micro-F1
intervals above are 0.19--0.25 wide, so modest real effects would go
undetected~\cite{card2020}.

\paragraph{Severity.} The rerun also scored harm severity (1--5). Its accuracy,
0.50, is below that of always predicting the modal score 1 (0.78). Its
quadratic-weighted $\kappa$, however, is 0.33 ($[0.06, 0.57]$) against 0 for the
constant: this is where the arm shows clear signal. The gold severities are
themselves inconsistent: four harm-labelled cases, such as ``Data warehouse'',
have severity 1. One unparseable prediction was scored as 0, outside the scale.

\section{What the evaluation compared}
\label{sec:audit}

\paragraph{Five arms, one prompt.} In the committed code, the prompt differs by
arm only for rules-static. The deterministic-first hybrid adds text only when a
deterministic model's output exists, and the runner never supplies one. The RAG
arm receives no retrieved context. Building each arm's prompt offline yields
two distinct texts: one for rules-static, and one shared by the other five.
Their spread (micro-F1 0.195--0.285, EMR 0.580--0.620) is therefore an estimate
of replicate noise. Rules-static's micro-F1 lead over the best of them (0.106)
is about the size of that spread (0.090). The recorded table predates the
earliest committed runner (2025-11-24, in a private predecessor repository),
which is built the same way, so this finding is strong but not conclusive. The
``rules'' are five generic compliance sentences, such as ``If a data breach
occurs, affected users must be notified within 72 hours.'' They are unrelated
to the trained models. Yet the original summary describes the arm as
``deterministic regex/keyword heuristics'', although an LLM produced every one
of its labels.

\paragraph{Scorer, prompt, and model changed together.} The earliest committed
runner mapped gold labels to the four groups but scored raw predicted labels.
The current runner maps both. Scored the earlier way, the 2026 outputs fall to
Jaccard 0.041, micro-F1 0.167, and nDCG 0.041: below the 2025 row on every
metric. The prompt changed too. A revision on 2025-11-29 added a root-cause
allowlist and told the model it may put root causes in the same field. The
model is unrecorded in both runs, and the provider default in 2025 was already
the model named in the 2026 README. The 2025$\to$2026 differences are
therefore a confound of scorer, prompt, possibly model, and one trial against a
three-trial mean. They are not evidence of model drift. Hosted-model drift is
real~\cite{chen2023,pozzobon2023}, but this artifact cannot measure it.

\paragraph{An LLM-generated, mostly off-topic gold set.} The summary records
that the v3 gold set was generated with Gemini~2.5~Pro. Arms from the same
family may share its errors. That is analogous to the self-preference of LLM
judges~\cite{panickssery2024,zheng2023}, and it would favour the arms, which
makes the constants' results more telling. Of the 50 cases, 47 are general
Wikipedia articles; by title, none reports a specific privacy incident. The
labelled cases mix on-topic articles (National Security Agency, Browser
security, Profiling) with off-topic ones (the Railway Regulation Act 1844, Data
warehouse). The JSON lists every title.

No case has the \texttt{harms} field the scorer reads first, so it silently
falls back to \texttt{root\_causes}. A code comment says that field holds
technical causes, but in this set it holds only Solove subtypes, so the
fallback read the right labels under the wrong name. The real contract failure
is on the prediction side. The 2026 prompt invites root causes into the field
the scorer treats as harms, which produced the eight out-of-taxonomy labels.

A second, 96-case set has labels drafted by \texttt{gemini-2.5-flash} (92 of
96) and is marked as reviewed in July 2026. Its review timestamps are bulk
writes (median gap 1.4\,ms), so they cannot show per-case review. It has no
empty-gold cases, which makes it the natural next benchmark, but no live
comparison on it survives.

\paragraph{Dry runs stored as experiments.} Fourteen ``comprehensive'' run
directories, the largest spanning 320 configurations, all name
\texttt{golden\_cases.jsonl} (the file that later became the 96-case set), not
the v3 set, so none bears on Table~\ref{tab:main}. Seven are flagged \texttt{dry\_run}. The
five that report per-arm metrics are all dry runs, and within four of them
every arm has an identical mean micro-F1; the fifth has two values. The
placeholder score, 0.417, is higher than the headline arm's 0.391, on a
different gold set. The seven live runs committed no per-arm metrics.
Pairwise-test files show Bonferroni thresholds and Cliff's~$\delta$ in three
directories. One file is truncated and does not parse, and 15 of the 43
parseable entries are placeholders ($p = 1$, zero difference). A related
failure path: on most API errors the LLM client returns a truncated echo of the
prompt. That parses to an empty prediction, which earns EMR credit on 70\% of
cases. This is an unexamined alternative explanation for part of the high 2025
EMRs; the rerun contains one such failure.

\paragraph{Heuristic labels in, heuristic rules out.} The labeller enables
GLiNER2 only when no LLM provider is configured, and it silently swallows a
GLiNER2 load failure, falling through to keyword scoring. In the released
33{,}472-row corpus, every label's recorded source is that keyword fallback.
61.4\% of labels are \texttt{unknown}, and 6 of the 21 other classes (e.g.\
\texttt{regulatory\_violation}) are not Solove categories. The models were
trained on seven features: a log count of affected individuals and six keyword
flags. So they approximated one keyword heuristic from another, which cannot
answer the study's first question, and they did so poorly:
\begin{itemize}
  \item Macro-F1 was 0.14 (decision tree), 0.10 (EBM), 0.20 (sparse linear),
        and 0.04 (Bayes net).
  \item The Bayes net predicted one class for every test row; that class makes
        up 52.4\% of its test split but only 2.1\% of the release. The release
        is therefore not the training distribution, and the training corpus's
        recorded hash matches no archived copy we could find.
  \item The accuracies reduce to 43/82, 191/328, 289/492, and 28/45. A test
        split shared by all four models would need at least 14{,}760 rows,
        more than a 20\% split of any corpus the repository documents, so the
        splits almost certainly differed.
  \item One of 16 per-category rule-list classifiers trained.
  \item The decision tree's 20 exported ``heuristics'' all report support
        0.000, yet its first rule (``no keyword present'') covers 61.4\% of
        release rows. The training corpus is unavailable, but zero support
        for that rule is implausible and most likely an exporter bug.
\end{itemize}

\section{Lessons}
\label{sec:lessons}

\begin{enumerate}[label=\textbf{L\arabic*}]
  \item \textbf{Pre-specify constant baselines and score them with the same
        scorer.} Use empty, all-labels in a fixed order, and the modal score for
        ordinal targets. Report oracle constants separately, re-selected inside
        the bootstrap. A metric on which no arm beats the pre-specified
        constants supports no claim.
  \item \textbf{Disclose empty-set conventions and the ceiling they imply.}
        Report the gold empty share, and score the non-empty subset separately.
  \item \textbf{Field names lie; pin the prompt--scorer contract.} Validate
        predictions and gold labels against one schema, and fail on missing
        fields or out-of-schema labels instead of falling back.
  \item \textbf{Diff what each arm actually sends.} Arms that differ only in
        name measure replicate noise. That is useful, but only if you know it.
  \item \textbf{Make placeholders unable to pass for results.} Refuse to
        aggregate or test arm-invariant outputs, keep dry runs out of results
        directories, and fail on client errors instead of echoing the prompt.
  \item \textbf{Keep the gold set independent and its review auditable.} Do not
        generate it with the model family under test, and record who reviewed
        each case and when.
  \item \textbf{Measure the label pipeline you ran, not the one you designed.}
        Log the share of labels from each source, and never swallow a model-load
        failure~\cite{ratner2017,pangakis2023}.
  \item \textbf{Record everything a comparison depends on, per result file.}
        That means the model version, scorer revision, prompt hash, trial count,
        and per-case outputs.
\end{enumerate}
\noindent\textbf{A data-ethics lesson.} A pre-release audit found that the
corpus's social-media portions republished verbatim, re-identifiable posts in
which private individuals described legal, medical, and financial trouble.
That is the secondary-use harm the project studied, and a tested scrub removed
it before release. Scrub at collection, not at release.

These lessons echo known failures: hypothesis-only and trivial
baselines~\cite{gururangan2018,poliak2018}, shortcut learning~\cite{geirhos2020},
leakage~\cite{kapoor2023}, fragile benchmark rankings~\cite{dehghani2021},
evaluation failures across ML~\cite{liao2021}, and significance
testing~\cite{dror2018}. What is new with LLMs is how cheaply one model family
can produce each ingredient (gold labels, rules, predictions, and even the
fallback text on error), and so end up on both sides of a comparison.

\section{Limitations}

\begin{itemize}
  \item We made no model calls, so we cannot say how the arms would score on
        the 96-case set.
  \item The exact scorer and prompt behind the 2025 table are unknown.
  \item The constant-versus-2025 comparison is nonetheless robust:
        \begin{itemize}
          \item the gold file is byte-identical to its Nov-2025 version;
          \item the taxonomy that maps the gold set's subtypes to groups is
                unchanged since the earliest committed runner, and the entries
                added later match no gold label;
          \item constants emit group labels, so their scores do not depend on
                how predictions are mapped.
        \end{itemize}
  \item If anything, the 2025 arms were handicapped by unmapped predictions.
        Mapping alone raises the rerun's Jaccard from 0.041 to 0.105.
  \item The filtered sensitivity check is post hoc.
  \item Our judgement that the gold articles are off-topic rests on titles.
  \item Evidence from the private predecessor repository cannot be checked by
        readers.
\end{itemize}

\section{Conclusion}

Pre-specified constants match or beat every recorded arm on three of the four
reported metrics. On the fourth, the recorded lead has no per-case evidence,
and the rerun shows signal only after its out-of-taxonomy labels are removed.
The study's first question, whether heuristics can estimate harm, remains open,
because every training label was itself a keyword heuristic. The uncertainty
machinery (confidence intervals, corrected tests) ran only on dry runs against
a different gold set, and the headline table carried no uncertainty at all.
What the study offers is a compact, reproducible case of evaluation failure,
and nine concrete ways to avoid repeating it.

\section*{Artifact availability}
\label{sec:artifact}
The repository, including gold sets, recorded runs, trained models, and a
scrubbed corpus (release \texttt{v1.0.0}), is public at
\url{https://github.com/Savage-Fred/privacy-harm-heuristics}.
\texttt{paper/reanalysis.py} regenerates \texttt{paper/reanalysis.json}, the
source of every re-analysis number, offline. It asserts that it reproduces the
recorded 2026 rerun (set, ranking, and ordinal metrics) and the repository's
offline arm exactly.

\begin{thebibliography}{99}\small
\bibitem{card2020} D.~Card, P.~Henderson, U.~Khandelwal, R.~Jia, K.~Mahowald, and D.~Jurafsky. With little power comes great responsibility. In \emph{EMNLP}, 2020.
\bibitem{chen2023} L.~Chen, M.~Zaharia, and J.~Zou. How is ChatGPT's behavior changing over time? arXiv:2307.09009, 2023.
\bibitem{chow1968} C.~K. Chow and C.~N. Liu. Approximating discrete probability distributions with dependence trees. \emph{IEEE Trans.\ Information Theory}, 14(3):462--467, 1968.
\bibitem{citron2022} D.~K. Citron and D.~J. Solove. Privacy harms. \emph{Boston University Law Review}, 102:793, 2022.
\bibitem{dehghani2021} M.~Dehghani, Y.~Tay, A.~A. Gritsenko, Z.~Zhao, N.~Houlsby, F.~Diaz, D.~Metzler, and O.~Vinyals. The benchmark lottery. arXiv:2107.07002, 2021.
\bibitem{dror2018} R.~Dror, G.~Baumer, S.~Shlomov, and R.~Reichart. The hitchhiker's guide to testing statistical significance in natural language processing. In \emph{ACL}, 2018.
\bibitem{geirhos2020} R.~Geirhos, J.-H. Jacobsen, C.~Michaelis, R.~Zemel, W.~Brendel, M.~Bethge, and F.~A. Wichmann. Shortcut learning in deep neural networks. \emph{Nature Machine Intelligence}, 2:665--673, 2020.
\bibitem{gururangan2018} S.~Gururangan, S.~Swayamdipta, O.~Levy, R.~Schwartz, S.~R. Bowman, and N.~A. Smith. Annotation artifacts in natural language inference data. In \emph{NAACL-HLT}, 2018.
\bibitem{jarvelin2002} K.~J\"arvelin and J.~Kek\"al\"ainen. Cumulated gain-based evaluation of IR techniques. \emph{ACM Trans.\ Information Systems}, 20(4):422--446, 2002.
\bibitem{kapoor2023} S.~Kapoor and A.~Narayanan. Leakage and the reproducibility crisis in machine-learning-based science. \emph{Patterns}, 4(9):100804, 2023.
\bibitem{lee2024} H.-P. Lee, Y.-J. Yang, T.~S. von Davier, J.~Forlizzi, and S.~Das. Deepfakes, phrenology, surveillance, and more! A taxonomy of AI privacy risks. In \emph{CHI}, 2024.
\bibitem{letham2015} B.~Letham, C.~Rudin, T.~H. McCormick, and D.~Madigan. Interpretable classifiers using rules and Bayesian analysis: Building a better stroke prediction model. \emph{Annals of Applied Statistics}, 9(3):1350--1371, 2015.
\bibitem{liao2021} T.~Liao, R.~Taori, I.~D. Raji, and L.~Schmidt. Are we learning yet? A meta review of evaluation failures across machine learning. In \emph{NeurIPS Datasets and Benchmarks}, 2021.
\bibitem{lou2012} Y.~Lou, R.~Caruana, and J.~Gehrke. Intelligible models for classification and regression. In \emph{KDD}, 2012.
\bibitem{lou2013} Y.~Lou, R.~Caruana, J.~Gehrke, and G.~Hooker. Accurate intelligible models with pairwise interactions. In \emph{KDD}, 2013.
\bibitem{nori2019} H.~Nori, S.~Jenkins, P.~Koch, and R.~Caruana. InterpretML: A unified framework for machine learning interpretability. arXiv:1909.09223, 2019.
\bibitem{pangakis2023} N.~Pangakis, S.~Wolken, and N.~Fasching. Automated annotation with generative AI requires validation. arXiv preprint, 2023.
\bibitem{panickssery2024} A.~Panickssery, S.~R. Bowman, and S.~Feng. LLM evaluators recognize and favor their own generations. In \emph{NeurIPS}, 2024.
\bibitem{poliak2018} A.~Poliak, J.~Naradowsky, A.~Haldar, R.~Rudinger, and B.~Van~Durme. Hypothesis only baselines in natural language inference. In \emph{*SEM}, 2018.
\bibitem{pozzobon2023} L.~Pozzobon, B.~Ermis, P.~Lewis, and S.~Hooker. On the challenges of using black-box APIs for toxicity evaluation in research. In \emph{EMNLP}, 2023.
\bibitem{ratner2017} A.~Ratner, S.~H. Bach, H.~Ehrenberg, J.~Fries, S.~Wu, and C.~R\'e. Snorkel: Rapid training data creation with weak supervision. \emph{PVLDB}, 11(3):269--282, 2017.
\bibitem{solove2006} D.~J. Solove. A taxonomy of privacy. \emph{University of Pennsylvania Law Review}, 154(3):477--564, 2006.
\bibitem{zaratiana2024} U.~Zaratiana, N.~Tomeh, P.~Holat, and T.~Charnois. GLiNER: Generalist model for named entity recognition using bidirectional transformer. In \emph{NAACL}, 2024.
\bibitem{zaratiana2025} U.~Zaratiana et~al. GLiNER2: An efficient multi-task information extraction system with schema-driven interface. arXiv preprint, 2025.
\bibitem{zhang2014} M.-L. Zhang and Z.-H. Zhou. A review on multi-label learning algorithms. \emph{IEEE Trans.\ Knowledge and Data Engineering}, 26(8):1819--1837, 2014.
\bibitem{zheng2023} L.~Zheng et~al. Judging LLM-as-a-judge with MT-Bench and Chatbot Arena. In \emph{NeurIPS Datasets and Benchmarks}, 2023.
\end{thebibliography}

\end{document}
